\documentclass[DIN, pagenumber=false, fontsize=11pt, parskip=half]{scrartcl}

\usepackage[utf8]{inputenc}
\usepackage[T1]{fontenc}
\usepackage{template}
\usepackage{amsmath}

\setcounter{secnumdepth}{1} % number sections only; subsections (Exercises) unnumbered

\titlehead{\centering\small
    Worked solutions to selected exercises from\\
    G. Chrystal, \emph{Algebra: An Elementary Text-Book}, Part I (1886)\\
\rule{0.4\linewidth}{0.4pt}}

\title{}
\author{} % Oscar A. Carrera
\date{}

\begin{document}
\maketitle

\section{The Fundamental Laws and Processes of Algebra}

\subsection{Exercises I}

\begin{problems}[from=1,step=1]
    \problem Point out in what sense the usual arrangement of the multiplication of 365 by 492 is an instance of the law of distribution.

    \answer The usual arrangement is
    \[
        \begin{array}{r@{\quad}r@{\qquad}l}
                   &  365 \\
            \times &  492 \\ \cline{1-2}
                   &  730              & \leftarrow 365 \times 2 \\
                   &  3285\phantom{0}  & \leftarrow 365 \times 90 \\
                 +  &  1460\phantom{00} & \leftarrow 365 \times 400 \\ \cline{1-2}
                   &  179580
        \end{array}
    \]
    Decimal notation already gives the multiplier as a sum, one term per digit: $492 = 400 + 90 + 2$. The arrangement therefore computes
    \[
        365 \times 492 = 365 \times (2 + 90 + 400) = 365 \times 2 + 365 \times 90 + 365 \times 400.
    \]
    This is the \emph{law of distribution}, $a(b + c + d) = ab + ac + ad$, with $a = 365$, $b = 2$, $c = 90$, $d = 400$. 

    The shifting of the second and third rows comes from writing $365 \times 90 = (365 \times 9) \times 10$, an instance of association, together with the decimal place-shift. 
    

    \problem I have a multiplying machine, but the most it can do at one time is to multiply a number of 10 digits by another number of 10 digits. Explain how I can use my machine to multiply 13693456783231 by 46381239245932.

    \answer First we cut the multiplier and multiplicand into digits we can enter into the machine. For decimal notation, we can do this preparation by splitting them into a sum of two units.
    \begin{align*}
        13693456783231 & = 1369345 \times 10^{7} + 6783231\\ 
        46381239245932 & = 4638123 \times 10^{7} + 9245932
    \end{align*}
    Now each number is of the form $x_2 \times 10^{7} + x_1$, with the two variables consisting of 7 digits each. 

    Now, giving the multiplicand and multiplier the assignments $a = 1369345$, $b = 6783231$, $c = 4638123$ and $d = 9245932$ and by application of the \emph{law of distribution}
    \begin{align*}
        \MoveEqLeft 13693456783231 \times 46381239245932 \\ 
            & = (a \times 10^7 + b) \times (c \times 10^7 + d)\\ 
            & = a \times 10^7 \times c \times 10^7 + ad \times 10^7 + bc \times 10^7 + bd\\
            & = ac \times 10^{14} + (ad + bc) \times 10^7 + bd
    \end{align*}
    Now we can use our multiplying machine to perform the four multiplications, namely $ac, ad, bc$ and $bd$. The remaining operations, multiplications by a power of 10 and simple addition, can be done by hand.


    \problem To divide 5004 by 12 is the same as to divide 5004 by 3, and then divide the quotient thus obtained by 4. Of what law of algebra is this an instance?

    \answer This is an instance of the \textit{law of association}. 
    Since $12 = 3 \times 4$, we have the statement $5004 \div (3 \times 4) = 5004 \div 3 \div 4$.
    Notice the removal of the bracket, and since $\div$ precedes the bracket, every sign within has been reversed.

    \problem If the remainder on dividing $N$ by $a$ be $R$, and the quotient $P$, and if we divide $P$ by $b$ and find a remainder $S$, show that the remainder on dividing $N$ by $ab$ will be $aS + R$. Illustrate with $5015 \div 12$.

    \answer Let $Q$ be the quotient on dividing $P$ by $b$.
    Then the two divisions are the identities
    \[
        \begin{array}{lcr}
            N = aP + R, & P = bQ + S, & (R < a \text{ and } S < b)
        \end{array}
    \]
    Substituting the second in the first, and applying the laws of distribution and association.
    \[ 
        N = a(bQ+S)+R = abQ+aS+ R = ab \times Q + (aS+R)
    \]
    Thus $N$ is a multiple of $ab$ together with $aS+R$.
    Note that since $R \leq a-1$ and $S \leq b-1$,
    \[
        aS+R \leq a(b-1) + (a-1) = ab-1 < ab.
    \]
    Therefore on dividing $N$ by $ab$ the remainder is $aS+R$.

    For illustration, take $N=5015$, $a=3$ and $b=4$.
    Dividing 5015 by 3 gives $P=1671$, $R=2$.
    Dividing 1671 by 4 gives $Q=417$, $S=3$.
    The remainder on dividing 5015 by 12 is $aS + R = 3 \times 3 + 2 = 11$.
    And we see $5015 = 12 \times 417 + 11$.

    \problem Show how to multiply two numbers of 10 digits each so as to obtain merely the number of digits in the product, and the first three digits on the left of the product. Illustrate by finding the number of digits, and the first three left-hand digits in the following:
    \begin{parts}
        \item $3659893456789325678 \times 342973489379265$;
        \item $2^{64}$.
    \end{parts}

    \answer Cut each number after its first $k$ digits, so that
    \[
        x = a \times 10^{m} + r, \qquad y = c \times 10^{n} + s,
    \]
    where $a$ and $c$ have $k$ digits each and the tails satisfy $r < 10^{m}$, $s < 10^{n}$. By the law of distribution,
    \[
        xy = ac \times 10^{m+n} + (as \times 10^{m} + cr \times 10^{n} + rs).
    \]
    The bracketed part is the whole error made by neglecting the tails, and it is small: since $s < 10^{n}$, $r < 10^{m}$, and $rs < 10^{m+n}$,
    \[
        as \times 10^{m} + cr \times 10^{n} + rs < (a + c + 1) \times 10^{m+n}.
    \]
    Hence $xy$ lies between $ac \times 10^{m+n}$ and $(ac + a + c + 1) \times 10^{m+n}$. Now $ac$ has $2k - 1$ or $2k$ digits, whereas $a + c + 1$ has at most $k + 1$; so the correction $a + c + 1$ can disturb only the last $k + 1$ digits of $ac$, and the leading $k - 2$ or so digits, together with the number of digits, are left untouched. Thus the number of digits in $xy$ is the number of digits in $ac$ plus $m + n$, and the first digits of $xy$ are the first digits of $ac$ as long as the addition of $a + c + 1$ does not carry into the desired digits, if it does we can take $k$ to be larger.

    \textit{Illustration (a).} Take $k = 5$:
    \begin{align*}
        3659893456789325678 &= 36598 \times 10^{14} + 93456789325678, \\
        342973489379265 &= 34297 \times 10^{10} + 3489379265.
    \end{align*}
    Then $ac = 36598 \times 34297 = 1255201606$, and $a + c + 1 = 70896$, so
    \[
        1255201606 \times 10^{24} < xy < 1255272502 \times 10^{24}.
    \]
    Both bounds have $10 + 24 = 34$ digits and begin with $1255$. The product has therefore $34$ digits, and its first three digits are $125$.

    \textit{Illustration (b).} Since $2^{64} = (2^{32})^{2}$ and $2^{32} = 4294967296$ has exactly ten digits, this is a product of two numbers of ten digits each. With $k = 5$,
    \[
        4294967296 = 42949 \times 10^{5} + 67296,
    \]
    so $ac = 42949^{2} = 1844616601$ and $a + c + 1 = 85899$, giving
    \[
        1844616601 \times 10^{10} < 2^{64} < 1844702500 \times 10^{10}.
    \]
    Both bounds have $10 + 10 = 20$ digits and begin with $1844$. Hence $2^{64}$ has $20$ digits, the first three being $184$.

    \problem Express in the simplest form---
    \[
        -(-(-(-(\ \dots\ (-1)\ \dots\ )))),
    \]
    \begin{parts}
        \item where there are $2n$ brackets;
        \item where there are $2n+1$ brackets; $n$ being any whole number whatever.
    \end{parts}
    
    \answer 
    Let $N$ be the number of brackets. 
    Each of the N brackets is preceded by a $-$ sign, and by the law of association, removing such a bracket reverses the sign of its content. Since the innermost bracket contains -1, the value is $(-1)^N \times (-1) = -(-1)^N$.
    \begin{parts}
    \item If $N = 2n$, then 
        \begin{align*}
            -(-1)^N &= -(-1)^{2n}\\ 
                    &= -((-1)^2)^n\\ 
                    &= -(1^n)\\ 
                    &= -1
        \end{align*}
    \item If $N = 2n+1$, then 
        \begin{align*}
            -(-1)^N &= -(-1)^{2n+1}\\ 
                    &= -[(-1)^{2n}\times(-1)]\\ 
                    &= -[(+1) \times (-1)]\\ 
                    &= -(-1)\\ 
                    &= +1
        \end{align*}
    \end{parts}


    \problem Simplify and condense as much as possible---
    \[
        2a - \{3a - [a - (b - a)]\}.
    \]

    \problem Simplify---
    \begin{parts}
        \item $3\{4 - 5[6 - 7(8 - 9 \cdot \overline{10 - 11})]\}$;
        \item $\frac{1}{3}\{\frac{1}{4} - \frac{1}{5}[\frac{1}{6} - \frac{1}{7}(\frac{1}{8} - \frac{1}{9} \cdot \overline{\frac{1}{10} - \frac{1}{11}})]\}$.
    \end{parts}

    \problem Simplify---
    \[
        1 - (2 - (3 - (4 - \dots (9 - (10 - 11))\ \dots\ ))).
    \]

    \problem Distribute the following products:
    \begin{parts}
        \item $(a+b) \times (a+b)$;
        \item $(a-b) \times (a+b)$;
        \item $(3a - 6b) \times (3a + 6b)$;
        \item $(\tfrac{1}{3}a - \tfrac{1}{6}b) \times (\tfrac{1}{3}a + \tfrac{1}{6}b)$.
    \end{parts}

    \problem Simplify, by expanding and condensing as much as possible---
    \begin{align*}
        &\{(m+1)a + (n+1)b\}\{(m-1)a + (n-1)b\} \\
        {}+{} &\{(m+1)a - (n+1)b\}\{(m-1)a - (n-1)b\}.
    \end{align*}

    \problem Simplify---
    \[
        \left(x + \frac{1}{x}\right)\left(y + \frac{1}{y}\right) + \left(x - \frac{1}{x}\right)\left(y - \frac{1}{y}\right).
    \]

    \problem Simplify---
    \[
        \left(x + \frac{1}{x}\right)\left(y + \frac{1}{y}\right)\left(z + \frac{1}{z}\right) - \left(x - \frac{1}{x}\right)\left(y - \frac{1}{y}\right)\left(z - \frac{1}{z}\right).
    \]

    \problem Expand and condense as much as possible---
    \[
        (\tfrac{1}{2}x + \tfrac{1}{4}y + \tfrac{1}{8}z)(\tfrac{1}{2}x - \tfrac{1}{4}y + \tfrac{1}{8}z).
    \]
\end{problems}

\end{document}
